\documentclass[11pt,letterpaper]{article}

% ---------- Paquetes ----------
\usepackage[utf8]{inputenc}
\usepackage[T1]{fontenc}
\usepackage[spanish,es-tabla]{babel}
\usepackage[margin=2.5cm]{geometry}
\usepackage{graphicx}
\usepackage{booktabs}
\usepackage{longtable}
\usepackage{amsmath}
\usepackage{amssymb}
\usepackage{listings}
\usepackage{xcolor}
\usepackage{caption}
\usepackage{hyperref}
\usepackage{enumitem}
\usepackage{float}

\hypersetup{
    colorlinks=true,
    linkcolor=blue!50!black,
    citecolor=blue!50!black,
    urlcolor=blue!50!black,
    pdftitle={Informe CI/CD -- Equipo ACC},
    pdfauthor={Carlos Carpio, Cristhian Pacheco, Robinson Cando, Jeremy Alvarez}
}

\definecolor{codebg}{RGB}{245,245,245}
\definecolor{codecomment}{RGB}{100,100,100}
\definecolor{codekeyword}{RGB}{0,0,180}
\definecolor{codestring}{RGB}{160,30,30}

\lstdefinestyle{codigo}{
    backgroundcolor=\color{codebg},
    commentstyle=\color{codecomment}\itshape,
    keywordstyle=\color{codekeyword}\bfseries,
    stringstyle=\color{codestring},
    basicstyle=\footnotesize\ttfamily,
    breakatwhitespace=false,
    breaklines=true,
    captionpos=b,
    keepspaces=true,
    numbers=left,
    numbersep=5pt,
    numberstyle=\tiny\color{codecomment},
    showspaces=false,
    showstringspaces=false,
    showtabs=false,
    tabsize=2,
    frame=single,
    rulecolor=\color{black!20}
}
\lstset{style=codigo}

\title{\textbf{Informe CI/CD}\\
\large Diagnóstico y corrección de fallas reales en el pipeline de integración continua\\
\large Equipo ACC --- Aplicaciones Distribuidas (ISR-701)}

\author{
Carlos Jos\'e Carpio Mendoza \and
Cristhian Daniel Pacheco C\'ardenas \and
Robinson Rodrigo Cando Moreno \and
Jeremy Alexis \'Alvarez P\'arraga\\[1em]
\normalsize Equipo ACC\\
\normalsize Facultad de Ciencias de la Computaci\'on\\
\normalsize Universidad T\'ecnica Estatal de Quevedo
}

\date{Agosto 2026}

\begin{document}

\maketitle

\begin{abstract}
\noindent Este informe documenta la depuración real del job \texttt{integration} del pipeline
de CI/CD (7 jobs), que fallaba de forma consistente en GitHub Actions pese a funcionar sin
problema en las máquinas de desarrollo del equipo. Se describen, en el orden en que se
descubrieron, tres causas raíz independientes -- todas confirmadas con evidencia real de logs,
nunca supuestas -- incluyendo los intentos previos que \emph{no} resolvieron el problema, la
corrección definitiva de cada una, y la verificación de punta a punta contra el stack real, hasta
dejar los 7 jobs en verde sobre el mismo commit. Se documentan también los bugs secundarios
encontrados en el propio proceso de verificar cada corrección, y las lecciones que deja el
patrón común a los tres problemas.
\end{abstract}

\tableofcontents
\newpage

\section{Contexto}

El equipo ACC mantiene un pipeline de integración continua en
\texttt{.github/workflows/ci-cd.yml}, compuesto por 7 jobs encadenados mediante \texttt{needs:}
para que cada etapa solo se ejecute si la anterior pasó:

\begin{table}[H]
\centering
\caption{Los 7 jobs del pipeline}
\begin{tabular}{@{}p{3.2cm}p{10.8cm}@{}}
\toprule
\textbf{Job} & \textbf{Qué hace} \\
\midrule
\texttt{lint} & Análisis estático de todos los subproyectos (backend, web, móvil). \\
\texttt{test-backend} & Pruebas unitarias e integración de los microservicios Java, Node.js y
Python. \\
\texttt{test-web} & Pruebas unitarias y de contrato (Pact) de \texttt{apps/web}. \\
\texttt{test-mobile} & Pruebas unitarias de \texttt{apps/mobile} (Kotlin/Jetpack Compose). \\
\texttt{build-images} & Matriz de 7 destinos: construye y publica en GHCR la imagen Docker de
cada microservicio y de \texttt{apps/web}. Solo publica en push, no en pull request. \\
\texttt{build-mobile-apk} & Compila el APK de depuración de \texttt{apps/mobile} y lo publica
como artefacto del pipeline. \\
\texttt{integration} & Levanta el stack completo con \texttt{docker compose up -d --build} y
corre contra él: verificación del proveedor del contrato Pact y la suite completa de Playwright
E2E en 3 navegadores. \\
\bottomrule
\end{tabular}
\end{table}

Los primeros 6 jobs pasaban de forma consistente. El séptimo, \texttt{integration}, fallaba en
\emph{todas} las corridas sobre GitHub Actions, mientras que el mismo
\texttt{docker compose up -d --build} que ese job ejecuta funcionaba sin problema en las máquinas
de desarrollo del equipo. Esa asimetría -- funciona en local, falla siempre en el runner -- fue
la pista central: el problema no era una falla puntual ni un error de sintaxis, sino una
diferencia estructural entre ambos entornos que el equipo tuvo que identificar con precisión, no
adivinar.

\subsection{Metodología de diagnóstico}

El primer obstáculo fue de \emph{visibilidad}, no de código: la API pública de GitHub
(\texttt{checks-runs/.../annotations}) solo devuelve mensajes genéricos del tipo
\texttt{"Process completed with exit code 1"}, sin el log real del comando que falló, y
descargar el log completo de un run sin autenticación devuelve \texttt{403 Must have admin
rights}. Con solo esa API, cada intento de arreglo era en el fondo una conjetura educada sobre
qué pudo haber fallado.

El equipo instaló y autenticó la CLI oficial de GitHub (\texttt{gh}), que sí da acceso al log
completo de cada paso individual del job (\texttt{gh run view <id> --log-failed}). Este cambio
de herramienta fue lo que permitió pasar de "adivinar la causa y esperar a ver si el siguiente
run pasa" a "leer el error real, confirmarlo, y recién entonces corregirlo" -- el patrón que
se repite en los tres problemas de este informe. En Windows, la instalación de \texttt{gh} vía
\texttt{winget} presentó además una complicación menor: el \texttt{PATH} actualizado no era
visible para sesiones de terminal ya abiertas, lo que obligó a invocar el binario por su ruta
absoluta hasta abrir una sesión nueva.

\section{Problema 1: candado circular en el healthcheck de CockroachDB}

\subsection{Síntoma inicial}
El paso \texttt{docker compose up -d --build} fallaba con
\texttt{dependency failed to start: container roach1 is unhealthy}, después de aproximadamente
70 segundos de espera.

\subsection{Primer intento: aumentar el presupuesto de tiempo (insuficiente)}
La hipótesis inicial, razonable dado que un runner compartido de GitHub Actions tiene menos CPU
disponible que una máquina de desarrollo dedicada, fue que CockroachDB simplemente necesitaba más
tiempo para formar el clúster bajo contención de recursos. Se subió el presupuesto del
healthcheck:

\begin{lstlisting}[language=yaml,caption={Primer intento: mas tiempo, no resolvio el problema}]
healthcheck:
  test: ["CMD", "curl", "-f", "http://localhost:8080/health?ready=1"]
  interval: 10s
  timeout: 5s
  retries: 15        # antes: 5
  start_period: 60s   # antes: 20s
\end{lstlisting}

El resultado del siguiente run real: la falla \emph{seguía ocurriendo}, ahora después de
aproximadamente 3 minutos y 21 segundos en vez de 70 segundos. Aumentar el tiempo no cambió el
resultado, solo lo retrasó -- la señal clara de que la causa no era de tiempo, sino estructural.

\subsection{Diagnóstico con logs reales}
Como ningún log de CockroachDB se había capturado nunca (el propio \texttt{docker compose up}
fallaba y abortaba antes de llegar a un paso posterior que sí volcaba logs), se modificó el paso
del workflow para volcar los logs de los 3 nodos ante cualquier falla:

\begin{lstlisting}[language=bash,caption={Paso de diagnostico agregado al workflow}]
docker compose up -d --build || {
  echo "FALLO 'docker compose up' -- logs de los 3 nodos de CockroachDB:"
  docker compose logs roach1 roach2 roach3
  exit 1
}
\end{lstlisting}

El siguiente run reveló, por fin, el mensaje real de CockroachDB en los 3 nodos:
\texttt{"The server appears to be unable to contact the other nodes in the cluster... Node will
now attempt to join a running cluster, or wait for \texttt{cockroach init}"}.

\subsection{Causa raíz real}
El healthcheck de los 3 nodos usaba \texttt{/health?ready=1} -- un endpoint de
\emph{readiness} que en CockroachDB solo responde \texttt{OK} una vez que el nodo ya pertenece a
un clúster \textbf{inicializado}. El servicio \texttt{db-init}, que es quien ejecuta
\texttt{cockroach init} para inicializar ese clúster, tenía como condición
\texttt{depends\_on: roach1/2/3: service\_healthy} -- es decir, esperaba a que los nodos
estuvieran "listos" antes de poder inicializarlos.

Candado circular real: nadie queda listo porque nadie los inicializa, y nadie los inicializa
porque nadie queda listo. Ningún valor de \texttt{retries}/\texttt{start\_period}, por grande que
fuera, podía resolver esto -- el candado nunca se abre por sí solo, sin importar cuánto tiempo se
le dé.

En las máquinas de desarrollo del equipo esto nunca se manifestó porque los volúmenes de Docker
ya traían un clúster inicializado de sesiones de trabajo anteriores: los nodos arrancaban,
encontraban su directorio de datos ya inicializado, y pasaban el healthcheck de inmediato. El
runner de GitHub Actions, en cambio, arranca siempre desde un volumen completamente vacío, así
que el candado se activa sin excepción, cien por ciento de las veces.

\subsection{Solución definitiva}
\begin{itemize}[leftmargin=1.2em]
  \item Los 3 healthchecks se cambiaron a \texttt{/health} (sin \texttt{?ready=1}) -- verifica
  solo que el proceso está vivo (\emph{liveness}), no que el clúster ya esté inicializado. Esto
  desbloquea a \texttt{db-init} para que pueda arrancar y ejecutar \texttt{cockroach init}.
  \item Se agregó un bucle de reintentos alrededor de \texttt{cockroach init} (hasta 30 intentos,
  2 s entre cada uno) y una segunda espera activa a que el clúster acepte conexiones SQL antes de
  correr los scripts de esquema, para cerrar una carrera residual pequeña entre "el nodo terminó
  de inicializarse" y "el nodo ya acepta SQL real".
\end{itemize}

\begin{lstlisting}[language=bash,caption={Extracto del entrypoint corregido de db-init}]
for i in $(seq 1 30); do
  out=$(cockroach init --insecure --host=roach1 2>&1) && break
  echo "$out" | grep -qi 'already been initialized' && break
  echo "db-init: intento $i fallo, reintentando en 2s: $out"
  sleep 2
done &&
for i in $(seq 1 30); do
  cockroach sql --insecure --host=roach1 -e 'SELECT 1;' >/dev/null 2>&1 && break
  echo "db-init: el cluster aun no acepta SQL (intento $i), reintentando en 2s"
  sleep 2
done &&
cockroach sql --insecure --host=roach1 -f /scripts/init_db.sql && ...
\end{lstlisting}

\subsection{Verificación}
Antes de subir el cambio se reprodujo el escenario de un entorno limpio en local, con una red y
volúmenes de Docker completamente nuevos (sin reutilizar los datos ya inicializados del equipo),
confirmando de punta a punta que el clúster se inicializaba correctamente sin intervención
manual. La siguiente corrida real en GitHub Actions confirmó
\texttt{roach1/roach2/roach3: Healthy} y el job \texttt{integration} avanzó más allá de este
punto por primera vez en todo el proceso de depuración.

\section{Problema 2: datos de demo faltantes en el contrato Pact}

\subsection{Síntoma}
Resuelto el problema anterior, el job avanzó considerablemente más lejos -- hasta la
verificación del lado proveedor del contrato Pact, penúltimo paso del job -- y ahí falló con:

\begin{lstlisting}[caption={Fallo real de la verificacion del contrato}]
1) Verifying a pact between soporte-web and ticket-service - una peticion GET
   de la lista de tickets, sin filtros, con un JWT valido de ADMIN has a
   matching body

   1.1) body: $.data Variant at index 0 ({"ticketId":"c8e57689-...",
        "status":"RESUELTO", ...}) was not found in the actual list
\end{lstlisting}

El ticket que el contrato esperaba encontrar en la respuesta real de \texttt{ticket-service}
simplemente no existía en la base de datos.

\subsection{Investigación}
El método \texttt{@State} de la prueba (\texttt{TicketServiceProviderPactTest}) era, hasta ese
momento, un no-op a propósito, con un comentario explícito que decía que el stack "ya tiene datos
de demo con tickets visibles para ADMIN", citando un script,
\texttt{resultados/rebalance\_demo\_tickets.sql}, como la fuente de esos datos. Se buscó ese
archivo en el repositorio completo, incluyendo todo el historial de commits, y se confirmó que
\textbf{nunca se llegó a versionar en git} -- no existe ni existió jamás en ningún commit. Solo
había existido, de forma manual y nunca documentada como tal, en el volumen de Docker de la
persona que lo escribió originalmente en su propia máquina.

Es exactamente el mismo patrón de fondo que el Problema 1: la máquina de desarrollo
"funcionaba" solo porque conservaba un estado acumulado que un entorno realmente limpio nunca
podría reproducir por su cuenta.

\subsection{Solución}
Se reescribió el método \texttt{@State} para que inserte, por SQL directo, exactamente los dos
tickets que el archivo de contrato (\texttt{pacts/soporte-web-ticket-service.json}) necesita --
es precisamente para eso que existe el mecanismo \texttt{@State} de Pact: que la propia prueba
garantice su precondición de forma explícita y reproducible, en vez de depender de un dato
externo invisible que nadie versionó.

\begin{lstlisting}[language=Java,caption={Extracto del @State corregido}]
@State("el rol ADMIN tiene al menos un ticket visible")
void elRolAdminTieneAlMenosUnTicketVisible() throws Exception {
    try (Connection conn = DriverManager.getConnection(
            "jdbc:postgresql://localhost:26257/ticket_db?sslmode=disable",
            "root", "")) {
        // ... inserta por SQL directo exactamente los tickets que el
        // archivo de contrato espera encontrar.
    }
}
\end{lstlisting}

\subsection{Bug adicional encontrado durante la verificación: \texttt{UPSERT} vs. índice único}
Al probar esta corrección contra el stack real (no simulada) se encontró un segundo bug, más
pequeño pero real: el primer intento usó \texttt{UPSERT INTO} para insertar los tickets, y falló
con \texttt{duplicate key value violates unique constraint "tickets\_id\_key"}.

La causa: \texttt{UPSERT} en CockroachDB solo resuelve conflictos contra la \emph{primary key}
compuesta de la tabla (\texttt{created\_at}, \texttt{id}); la tabla \texttt{tickets} tiene además
un índice único \emph{aparte} sobre \texttt{id} en solitario (\texttt{tickets\_id\_key}, agregado
para que las búsquedas por id no tengan que escanear las 4 particiones). Como esos tickets de
demo ya existían de pruebas manuales anteriores en el mismo stack, con un valor distinto de
\texttt{created\_at}, el \texttt{UPSERT} no los reconoció como el mismo registro y chocó contra
ese índice en vez de actualizar la fila.

\begin{lstlisting}[language=SQL,caption={Reemplazo de UPSERT por DELETE + INSERT}]
-- Antes (fallaba si el ticket ya existia con otro created_at):
UPSERT INTO tickets (created_at, id, ...) VALUES (?, ?, ...)

-- Despues (idempotente sin importar el estado previo de la tabla):
DELETE FROM tickets WHERE id = ?;
INSERT INTO tickets (created_at, id, ...) VALUES (?, ?, ...);
\end{lstlisting}

\subsection{Verificación}
Se corrió la prueba dos veces seguidas contra el stack real con Maven real (no una ejecución
simulada): la primera corrida expuso el bug del \texttt{UPSERT} en vivo, con el error real de
CockroachDB; tras aplicar el fix, ambas corridas pasaron 2/2 de forma consistente, confirmando
la idempotencia real del mecanismo sin importar si los tickets ya existían de una corrida
anterior.

\section{Problema 3: Playwright solo instalaba Chromium}

\subsection{Síntoma}
Con los dos problemas anteriores resueltos, el job llegó por primera vez hasta las pruebas E2E de
Playwright. Los 3 casos de Chromium pasaron:

\begin{lstlisting}[caption={Chromium paso, Firefox y WebKit fallaron de inmediato}]
  ok 1 [chromium] un ADMIN inicia sesion y ve la consola con tickets reales
  ok 2 [chromium] abrir un ticket muestra su detalle completo
  ok 3 [chromium] credenciales invalidas muestran un error real del backend

  1) [firefox] un ADMIN inicia sesion...
     Error: browserType.launch: Executable doesn't exist at
     /home/runner/.cache/ms-playwright/firefox-1538/firefox/firefox
\end{lstlisting}

Firefox y WebKit fallaban de inmediato, sin siquiera llegar a intentar abrir una página --
el navegador simplemente no estaba instalado en el runner.

\subsection{Causa raíz}
\texttt{playwright.config.ts} define 3 proyectos (\texttt{chromium}, \texttt{firefox},
\texttt{webkit} -- requisito de "Compatibilidad" de la rúbrica de Entrega 4: la web debe
funcionar en Chrome, Firefox y Safari, usando WebKit como el mismo motor de renderizado que usa
Safari). El paso de instalación del workflow, sin embargo, decía:

\begin{lstlisting}[language=bash,caption={Paso de instalacion original, restringido a un solo navegador}]
npx playwright install --with-deps chromium
\end{lstlisting}

El argumento \texttt{chromium} restringía la instalación a un solo navegador de los tres que la
configuración realmente usa -- una desconexión simple entre lo que el \texttt{.yml} del pipeline
instalaba y lo que el propio \texttt{playwright.config.ts} del proyecto exigía.

\subsection{Solución}
\begin{lstlisting}[language=bash,caption={Correccion: sin argumento, instala los navegadores que la config usa}]
npx playwright install --with-deps
\end{lstlisting}

\subsection{Verificación}
Se instalaron los 3 navegadores en local (confirmando la presencia real de los binarios de
\texttt{firefox} y \texttt{webkit} en la caché de Playwright) y se corrieron las 9 pruebas E2E
completas contra el stack real levantado, confirmando que ningún caso volvía a fallar con el
error de ejecutable faltante, en ningún navegador, antes de subir el cambio.


\section{Resultado}

Con los tres cambios integrados en commits sucesivos sobre la misma rama, el pipeline completo --
los 7 jobs, incluidos los 7 destinos de la matriz de \texttt{build-images} y
\texttt{build-mobile-apk} -- quedó en verde sobre un mismo commit, cumpliendo el requisito de la
rúbrica de que los 7 jobs pasen juntos en el último commit de la entrega.

\begin{table}[H]
\centering
\caption{Resumen de las tres causas raíz encontradas y corregidas}
\begin{tabular}{@{}p{2.6cm}p{4.6cm}p{4.4cm}p{2.6cm}@{}}
\toprule
\textbf{Problema} & \textbf{Causa raíz} & \textbf{Corrección} & \textbf{Intento previo fallido} \\
\midrule
Healthcheck de CockroachDB & Candado circular: \texttt{?ready=1} exige clúster ya inicializado,
pero quien inicializa espera a ese mismo healthcheck & \texttt{/health} (liveness) + reintentos
de \texttt{cockroach init} + espera activa a SQL & Aumentar \texttt{retries}/
\texttt{start\_period} -- solo retrasó la falla \\
\addlinespace
Contrato Pact & Datos de demo asumidos como ya existentes, en un script nunca versionado en git &
El propio \texttt{@State} inserta los datos por SQL (\texttt{DELETE}+\texttt{INSERT}) &
\texttt{UPSERT} -- chocó contra un índice único aparte \\
\addlinespace
Playwright E2E & Paso de instalación restringido a un solo navegador (\texttt{chromium}) &
Instalar los 3 navegadores que la configuración define & -- \\
\bottomrule
\end{tabular}
\end{table}

\section{Amenazas a la validez}

\begin{itemize}[leftmargin=1.2em]
  \item \textbf{Un solo entorno de runner observado.} Todas las corridas se hicieron sobre los
  runners \texttt{ubuntu-latest} compartidos de GitHub Actions; no se probó el pipeline en otro
  proveedor de CI, así que no se puede afirmar que estas sean las únicas causas posibles de falla
  en un entorno "limpio" en general, solo que lo fueron en el entorno real usado por el equipo.
  \item \textbf{El primer intento del Problema 1 no se descartó por análisis, sino por
  observación.} Aumentar el presupuesto de tiempo era una hipótesis razonable a priori; se supo
  que era insuficiente solo después de ver el resultado del siguiente run real, no por haber
  anticipado el candado circular de antemano. Esto es consistente con la metodología descrita
  (diagnosticar con logs reales antes de dar un fix por bueno), pero vale dejarlo explícito.
  \item \textbf{La verificación de idempotencia del Problema 2 se hizo con 2 corridas, no más.}
  Es suficiente para confirmar que el bug del \texttt{UPSERT} se corrigió, pero no constituye una
  prueba exhaustiva de que no exista alguna otra condición de carrera no observada en ese punto.
\end{itemize}

\section{Conclusiones}

Los tres problemas comparten un mismo patrón de fondo: algo que \emph{parecía} funcionar porque
la máquina de desarrollo del equipo conservaba un estado acumulado -- un volumen ya inicializado,
datos de demo cargados a mano, o simplemente un navegador ya instalado de una sesión anterior --
que un entorno realmente limpio, como el de un runner de GitHub Actions, nunca reproduce por su
cuenta. Ninguno de los tres era un problema de "CI/CD mal configurado" en abstracto: cada uno
tenía una causa raíz concreta y verificable con el log real del paso que fallaba, y cada
corrección -- incluidos los dos bugs secundarios encontrados en el propio proceso de verificar
las correcciones principales -- se probó de punta a punta contra el stack real, nunca simulada,
antes de darla por buena.

Una lección práctica adicional: el primer intento de arreglar el Problema 1 (aumentar tiempos de
espera) fue una hipótesis razonable que resultó incorrecta, y solo se supo que era incorrecta
observando el resultado real del siguiente run, no por haberla anticipado de antemano. Aumentar
un timeout es, en general, un cambio de bajo riesgo y fácil de justificar -- pero cuando la causa
real es estructural (un candado circular, en este caso), ningún valor de tiempo, por grande que
sea, lo resuelve. La señal de que había que seguir buscando no fue que el fix "no funcionara" en
abstracto, sino que el mismo síntoma reapareciera después de más tiempo en vez de desaparecer.

\end{document}
